\documentclass[10pt,a4paper]{article}
\usepackage{amsmath,amssymb,amsthm,graphicx}
\usepackage{amsmath,amssymb,amsthm}
\newtheorem{nttheorem}{Theorem}
\newtheorem{ntdefinition}[nttheorem]{Definition}
\newtheorem{ntlemma}[nttheorem]{Lemma}
\newtheorem{ntproposition}[nttheorem]{Proposition}
\newtheorem{ntclaim}[nttheorem]{Claim}
\title{Query-efficient distance amplification for non-adaptive locally decodable codes}
\author{Gil Cohen, Tal Yankovitz}
\date{doi:10.4230/LIPIcs.CCC.2021.1; CC BY4.0}
\begin{document}
\maketitle
\noindent\textit{Editorial note.} Gil Cohen and Tal Yankovitz. Rate Amplification and Query-Efficient Distance Amplification for Linear LCC and LDC. CCC2021, article1, pp1:1--1:57. DOI10.4230/LIPIcs.CCC.2021.1.\par\medskip
\noindent\textit{Editorial note.} License: https://creativecommons.org/licenses/by/4.0/\par\medskip
\noindent\textit{Editorial note.} Faithful editable transcription of the official CCC2021 CC BY4.0 publication, article1 pp1:1--1:57. This item collects the non-explicit distance-amplification main result Theorem72, complete Section6.1 construction and Section6.2 analysis including Proposition65 and full Claims66--71 proofs, Section6.2.1 complete proof and low-error/direct-instantiation discussion, Definition64, and supporting sampler Definition14/Theorem15 full author proof sketch/Lemma16 plus standard Theorem18. Theorem73 is retained as an original cited prerequisite and direct instantiation, not counted separately. All43 original bibliography entries retained. Original wording and indices, including Claim68 k3 and Claim69 v notin Bz, are preserved without mathematical correction. Author nonadaptivity, alphabets, constant rate/error, sampler degree conventions and all query/error dependencies retained. Explicit variants Sections6.3--6.6 are outside this item; the complete official original is preserved. No native TeX was found in the bounded source lookup; this is a transcription of the licensed official PDF. Root handles compilation and database delivery.\par\medskip
\noindent\textit{Allowed standard prerequisites (editorial).}\begin{itemize}
\item Chernoff bound and binary entropy bound. Original sampler proof sketch and Lemma16 retained; no recursion into standard probability proofs.
\item Gilbert-Varshamov codes. Original Theorem18 statement retained.
\item Asymptotically good locally decodable codes. Original Theorem73[26] statement retained as optional direct instantiation.
\end{itemize}\medskip
\setcounter{nttheorem}{71}
\begin{nttheorem}[Author Theorem 72]
Let $C$ be a block-length-$n$ $(q,\delta,1/5)$-LDC over alphabet $\Sigma$ having a constant rate. Let $C'$ be a family of asymptotically good $(q'_n,\delta',1/5)$-LDC, where $q'_n$ is the query complexity when the code from the family is taken with block length $n$. Then, there exists an asymptotically good LDC over $\Sigma$, with constant error, having block length $\Theta(n)$ and query complexity
\begin{equation}q_{\mathrm{new}}=O\left(q\cdot q'_{O(1/\delta)}\log(1/\delta)\log n\right).\tag{6.2}\end{equation}
\end{nttheorem}

\section*{Author background and Section 3.1: Samplers}
Recall that a function $C:\Sigma^k\to\Sigma^n$ is an error-correcting code with distance $\delta$ if for every distinct $x,y\in\Sigma^k$, $\operatorname{dist}(C(x),C(y))\geq\delta$, where $\operatorname{dist}$ is the relative Hamming distance. The rate of the code $C$ is given by $\rho=k/n$. When there is more than one possible message length, we consider a code family, which is a family of functions in which each function is a code, and there is one code per message length $k$. A code family is asymptotically good if both the rate and distance of every code in the family are uniformly bounded below by constants $\rho>0$ and $\delta>0$, respectively.
Our distance amplification procedure makes use of samplers. These are bipartite graphs with a certain pseudo-random property. Let $G=(L,R,E)$ be a bipartite graph. We say $G$ is left-regular if all nodes in $L$ have the same degree.
\setcounter{nttheorem}{13}
\begin{ntdefinition}[Author Definition 14, {\cite{ref9}}]
Let $0<\varepsilon,\delta<1$. A bipartite graph $G=(L,R,E)$ is an $(\varepsilon,\delta)$-sampler if for every subset $T\subseteq R$, for all but $\delta$-fraction of vertices $v\in L$ it holds that
\[\left|\frac{|\Gamma(v)\cap T|}{|\Gamma(v)|}-\mu(T)\right|\leq\varepsilon.\]
\end{ntdefinition}
We will be working with ``unbalanced'' samplers. These are samplers with $|L|\gg|R|$. The state-of-the-art constructions of these samplers rely on their connection to randomness seeded extractors. We refer the interested reader to the excellent survey by Goldreich~\cite{ref17} for more information. When working with samplers, it is rather typical that the bipartite graph is left-regular, that is, the degree of all vertices in $L$ is the same. A small additional technical property we need is that the degree of every vertex in $R$ is close to the average right-degree. We make use of the following theorem which gives (non-explicit) samplers with near-optimal parameters having the above properties with respect to the degrees. We give a proof sketch for completeness.
\begin{nttheorem}[Author Theorem 15]
There exists a universal constant $c_{\mathrm{samp}}\geq1$ such that the following holds. For all integers $\ell,r$ and all $\varepsilon>0$, $1/2>\delta>0$ for which $\ell\geq r/(\delta\log(1/\delta))$, there exists a left-regular $(\varepsilon,\delta)$-sampler $G=([\ell],[r],E)$ with left-degree $d=c_{\mathrm{samp}}\cdot\log(1/\delta)/\varepsilon^2$. Moreover, provided that $\log r<1/(\delta\varepsilon^2)$, every right vertex has degree in $[D/2,2D]$ where $D=\ell d/r$ is the average right degree.
\end{nttheorem}
For the proof we need the following well-known lemma.
\begin{ntlemma}[Author Lemma 16]
For every integers $1\leq k\leq n$ with $k/n=\delta\leq1/2$ it holds that
\[\sum_{i=0}^k\binom ni\leq2^{H(\delta)n},\]
where $H(x)=-x\log(x)-(1-x)\log(1-x)$ is the binary entropy function.
\end{ntlemma}
\begin{proof}[Author proof sketch for Theorem 15]
The proof is via the probabilistic method, where for every left vertex we choose $d$ neighbours independently and uniformly at random, and independently across all left vertices (note that in the above we allow for parallel edges, but if that troubles the reader, that can be avoided as well in the regime of interest $d\ll r$ by arguing that the probability of a right neighbor to be selected more than once is small. In any case, our distance amplification procedure works just as well with parallel edges). Fix $T\subseteq[r]$. For $v\in[\ell]$ let $F_v$ be the indicator random variables that is $1$ if and only if $\bigl||\Gamma(v)\cap T|/d-\mu(T)\bigr|>\varepsilon$. By the Chernoff bound, $\Pr[F_v]\leq e^{-\Omega(\varepsilon^2d)}$. Fix $S\subseteq[\ell]$ with $|S|=\delta\ell$. The probability that for all vertices $v\in S$ it holds that $F_v=1$ is bounded above by $e^{-\Omega(\varepsilon^2d\cdot\delta\ell)}$. By taking the union bound over all $S$ and $T$, we get that except with probability
\begin{equation}2^r\binom\ell{\delta\ell}e^{-\Omega(\varepsilon^2d\delta\ell)}\leq2^{r+H(\delta)\ell-c\varepsilon^2d\delta\ell},\tag{3.1}\end{equation}
the sampled graph is an $(\varepsilon,\delta)$-sampler. Note that the last inequality follows by Lemma 16, where $c>0$ is some constant. By taking $c_{\mathrm{samp}}\geq5/c$, one can verify (using that $H(x)\leq2x\log(1/x)$ for all $x\leq1/2$) that the right hand side in Equation (3.1) is bounded by $1/4$.
As for the moreover part, again, by the Chernoff bound, the probability that there exists a right vertex which has degree outside $[D/2,2D]$ is bounded above by $re^{-\Omega(\ell d/r)}$, and this is bounded by $1/4$ by our choice of parameters and by taking $c_{\mathrm{samp}}$ large enough.
\end{proof}
\setcounter{nttheorem}{17}
\begin{nttheorem}[Author Theorem 18: The Gilbert--Varshamov bound]
Let $\Sigma$ be a set of size $|\Sigma|=q$. For every $n\in\mathbb N$, and $0\leq\delta\leq1-1/q$ there exists a code of block-length $n$ over $\Sigma$, with distance at least $\delta$ and rate $r\geq1-H_q(\delta)$. Furthermore, if $q$ is a prime power and $\Sigma=\mathbb F_q$, there exists a linear code over $\Sigma$ with rate $r\geq1-H_q(\delta)-g(n)$, where $g(n)=O(1/n)$.
\end{nttheorem}

\section*{6 Query-efficient distance amplification}
In this section we construct our query-efficient distance amplification procedure. We start by giving a somewhat more formal definition of locally decodable codes (compared to Definition 1) or, more precisely, a more formal definition of their non-adaptive counterparts. Recall that, informally, these are LDC in which the joint distribution of queries depends solely on the index one wishes to decode and is independent of the received word. By inspection, it is our understanding that the AEL distance amplification procedure also requires non-adaptivity.
\setcounter{nttheorem}{63}
\begin{ntdefinition}[Author Definition 64: Locally decodable codes]
Let $(C,Q,R)$ be a tuple of functions
\begin{align*}C&:\Sigma_{\mathrm{in}}^k\to\Sigma_{\mathrm{out}}^n,\\Q&:[k]\times\{0,1\}^r\to[n]^q,\\R&:[k]\times\Sigma_{\mathrm{out}}^q\times\{0,1\}^r\to\Sigma_{\mathrm{in}}.\end{align*}
Define $D:[k]\times\Sigma_{\mathrm{out}}^n\times\{0,1\}^r\to\Sigma_{\mathrm{in}}$ as follows. For $v\in[k]$, $y\in\Sigma_{\mathrm{out}}^n$, and $s\in\{0,1\}^r$, let
\[Q(v,s)=(u_1,\ldots,u_q),\qquad D(v,y,s)=R(v,y_{u_1},\ldots,y_{u_q},s).\]
The tuple $(C,Q,R)$ is called a $(q,\delta,\varepsilon)$-locally decodable code (or $(q,\delta,\varepsilon)$-LDC for short) if the following holds. For every $v\in[k]$, $x\in\Sigma_{\mathrm{in}}^k$, and $y\in\Sigma_{\mathrm{out}}^n$ for which $\operatorname{dist}(y,C(x))\leq\delta$, it holds that
\[\Pr_{s\sim U_r}[D(v,y,s)=x_v]\geq1-\varepsilon.\]
We call the function $C$ the encoding function, $Q$ the querying function, and $R$ the reconstruction function. The induced function $D$ is called the decoding function. The parameters $k,n$ are referred to as the message length and the block length, respectively. The sets $\Sigma_{\mathrm{in}},\Sigma_{\mathrm{out}}$ are called the input alphabet and output alphabet, respectively. We will be interested mostly in locally decodable codes in which $\Sigma_{\mathrm{in}}=\Sigma_{\mathrm{out}}$ in which case we refer to both as the alphabet of the code. The parameter $r$ is called the randomness complexity of the LDC. We say the LDC is explicit if all three functions $C,Q,R$ are polynomial-time computable. Note that then the decoding function $D$ is also polynomial-time computable.
\end{ntdefinition}
\subsection*{6.1 The distance amplification procedure}
In this section we present our query-efficient distance amplification procedure. We start by describing the building blocks we use and specify their parameters.
\paragraph{Building blocks}
\begin{itemize}
\item For $i=1,2$ let $(C_i,Q_i,R_i)$ be a $(q_i,\delta_i,\varepsilon_i)$-LDC with message length $k_i$ and block length $n_i$ over the same alphabet $\Sigma$. We denote the rate $k_i/n_i$ of $C_i$ by $\rho_i$.
\item Let $(C_3,Q_3,R_3)$ be a family of $(q_3(k_3),\delta_3(k_3),\varepsilon_3(k_3))$-LDC having rate $\rho_3(k_3)$ for message length $k_3$. The code $C_3$ is also over the alphabet $\Sigma$. We will always work with functions $q_3,\delta_3,\varepsilon_3,\rho_3$ that are monotone. More precisely, $q_3$ and $\rho_3$ are non-decreasing and $\delta_3,\varepsilon_3$ are non-increasing. We sometimes write $q_3,\delta_3,\varepsilon_3,\rho_3$ without mentioning explicitly the message length, and by that refer to the largest $k_3$ considered in the construction for $q_3,\delta_3$ and the smallest $k_3$ when considering $\varepsilon_3,\rho_3$. In any case, we assume (mostly for simplicity) that $\rho_3(k_3)\geq1/2$ for all $k_3$.
\item Set $\ell=n_1/k_2$. Let $G=(L,R,E)$ be a $(\delta_2/2,\delta_1)$-sampler with $|L|=\ell$ and $|R|=r$. Assume $G$ is left-regular with left-degree $d=n_2$. Assume further that every right-vertex $v$ of $G$ has degree $\deg(v)\in[D/2,2D]$, where $D$ is the average right degree $D=\ell d/r=n_1/(r\rho_2)$.
\end{itemize}
\paragraph{How to think of the parameters?}
We think of $C_1$ as the code whose distance $\delta_1$ we wish to amplify. Typically, the code $C_2$ has a much shorter message length $n_2\ll n_1$. In all applications in this paper we take $\delta_2$ to be either constant or slightly sub constant in $n_1$. The code $C_3$ has a larger block length than $C_2$ and, depending on the application, it has either a somewhat smaller or much smaller message length than $n_1$. We typically take $\delta_3\approx\delta_2$. The rates of all three codes is taken to be constant and even close to one. Note that we take $C_3$ to be a family of codes, whereas $C_1$ and $C_2$ are codes with predetermined message lengths. The reason is that the sampler $G$ is not necessarily right-regular, and in the construction, we associate codes from $C_3$ with the right vertices of $G$. Recall, though that the ratio of largest to smallest right-degree is bounded by $4$, so that is a minor technicality.
To describe the LDC that is composed of these building blocks, we need to specify the encoding function, querying function and reconstruction function. We start by describing the encoding function.
\paragraph{The encoding function}
Let $n=\sum_{v\in R}n_v$ where $n_v$ is the block length of the code from the family $C_3$ having message length $k_v=\deg(v)$. We define the function $C:\Sigma^{k_1}\to\Sigma^n$ as follows. Let $x\in\Sigma^{k_1}$.
\begin{enumerate}
\item Compute $y=C_1(x)\in\Sigma^{n_1}$.
\item Partition $y$ to $y=y^{(1)}\circ\cdots\circ y^{(\ell)}$ consecutive blocks, each of length $k_2$. Recall that, indeed, $n_1=\ell k_2$.
\item For every $u\in[\ell]$ compute $z^{(u)}=C_2(y^{(u)})\in\Sigma^{n_2}$.
\item For every $v\in[r]$ and $j\in[\deg(v)]$ let $(u,j')=\Gamma(v,j)\in[\ell]\times[n_2]$. Define the string $w^{(v)}\in\Sigma^{\deg(v)}=\Sigma^{k_v}$ as follows: for $j\in[\deg(v)]$, $(w^{(v)})_j=(z^{(u)})_{j'}$.
\item For every $v\in[r]$ compute $t^{(v)}=C_3(w^{(v)})\in\Sigma^{n_v}$.
\item The output of the encoding function on input $x$ is then defined by $C(x)=t^{(1)}\circ\cdots\circ t^{(r)}\in\Sigma^n$, where as usual we identify $R$ with $[r]$.
\end{enumerate}
By the construction of the encoding function, the message length and block length of the resulted code are $k_1$ and $n$, respectively. From here on we denote $k=k_1$.
\paragraph{The querying function}
We denote the randomness complexity of $C_1,C_2,C_3$ by $r_1,r_2,r_3$, respectively. The randomness complexity of the resulting querying function will be $r=r_1+r_2+r_3$, and the query complexity will be $q\leq q_1q_2q_3$, where $q_3$ is taken to be the maximum query complexity taken over all right vertices. We turn to define the querying function $Q:[k]\times\{0,1\}^r\to[n]^q$ as follows. On inputs $p\in[k]$, $s\in\{0,1\}^r$ we proceed as follows.
\begin{enumerate}
\item Partition $s=s_1\circ s_2\circ s_3$ where $|s_1|=r_1$, $|s_2|=r_2$, $|s_3|=r_3$.
\item Compute $(a_1,\ldots,a_{q_1})=Q_1(p,s_1)\in[n_1]^{q_1}$.
\item For $i=1,\ldots,q_1$
\begin{enumerate}
\item Set $u_i=\lceil a_i/k_2\rceil$ and $b_i=1+((a_i-1)\bmod k_2)$. Informally, $u_i$ is the ``bucket'' in which $a_i$ resides and $b_i$ is its location within the bucket. Note that we start the counting from $1$ rather than $0$, hence the slightly annoying addition and subtraction by $1$ in the definition of $b_i$.
\item Compute $(t_1^{(i)},\ldots,t_{q_2}^{(i)})=Q_2(b_i,s_2)\in[n_2]^{q_2}$.
\item For $j=1,\ldots,q_2$
\begin{enumerate}
\item Let $(v^{(i,j)},\widehat t_j^{(i)})=\Gamma(u_i,t_j^{(i)})\in[r]\times[k_{v^{(i,j)}}]$.
\item Compute $(e_1^{(i,j)},\ldots,e_{q_3}^{(i,j)})=Q_3(\widehat t_j^{(i)},s_3)\in[n_{v^{(i,j)}}]^{q_3}$.
\item As before, we endow the right vertices of the sampler in a fixed (arbitrary) order by identifying $R$ with $[r]$. For $h=1,\ldots,q_3$ set $c^{(i,j,h)}$ to be the absolute location of $e_h^{(i,j)}$ in the ordering of $R$. That is, $c^{(i,j,h)}=e_h^{(i,j)}+\sum_{v<v^{(i,j)}}n_v$.
\end{enumerate}
\end{enumerate}
\item The result is then given by $Q(p,s)=(c^{(i,j,h)})_{(i,j,h)\in[q_1]\times[q_2]\times[q_3]}$.
\end{enumerate}
Note that, indeed, the query complexity $q$ of the querying function defined above is at most $q_1q_2q_3$ where, recall, $q_3=q_3(2D)$. From here on we identify $[q]$ with $[q_1]\times[q_2]\times[q_3]$.
\paragraph{The reconstruction procedure}
We define the reconstruction procedure $R:[k]\times\Sigma^q\times\{0,1\}^r\to\Sigma$ as follows. On inputs $p\in[k]$, $\sigma=(\sigma^{(i,j,h)})_{(i,j,h)\in[q_1]\times[q_2]\times[q_3]}\in\Sigma^q$, and $s\in\{0,1\}^r$, we proceed as follows.
\begin{enumerate}
\item Partition $s=s_1\circ s_2\circ s_3$ where $|s_1|=r_1$, $|s_2|=r_2$, $|s_3|=r_3$ as in the querying function.
\item For $i=1,\ldots,q_1$
\begin{enumerate}
\item For $j=1,\ldots,q_2$
\begin{enumerate}
\item Denote $(z_1,\ldots,z_{q_3})=(\sigma^{(i,j,1)},\ldots,\sigma^{(i,j,q_3)})$.
\item Compute $y_j^{(i)}=R_3(\widehat t_j^{(i)},z_1,\ldots,z_{q_3},s_3)$, where $\widehat t_j^{(i)}=\widehat t_j^{(i)}(p,s)$ as defined in the querying function.
\end{enumerate}
\item Set $x_i=R_2(b_i,y_1^{(i)},\ldots,y_{q_2}^{(i)},s_2)$ where $b_i=b_i(p,s)$ as defined in the querying function.
\end{enumerate}
\item The output is then given by $R(p,\sigma,s)=R_1(p,x_1,\ldots,x_{q_1},s_1)$.
\end{enumerate}

\subsection*{6.2 Analysis}
In this section we analyze the LDC obtained above. We prove
\setcounter{nttheorem}{64}
\begin{ntproposition}[Author Proposition 65]
With the notation of the previous section, $C$ is a $(q,\delta,\varepsilon)$-LDC, where
\[q\leq q_1q_2q_3,\qquad\delta\geq\frac{\delta_2\delta_3}{16},\qquad\varepsilon\leq\varepsilon_1+(\varepsilon_2+\varepsilon_3)n.\]
Furthermore, $C$ has rate $\rho_1\rho_2\rho_3$, where $\rho_1,\rho_2$ are as defined in the building blocks paragraph, and per our convention set above, $\rho_3=\rho_3(D/2)$.
\end{ntproposition}
\paragraph{Remark regarding the distance}
Note that the distance $\delta$ of the resulted code $C$ is independent of $\delta_1$, the poor distance of $C_1$ we set out to amplify. This is the key feature of the AEL distance amplification procedure (which our variant above, of course, maintains). It is yet another instance of a general strategy in pseudo-randomness that combines objects in such a way that the resulted object enjoys the upsides of the different parts and avoid their shortcomings. The Zig-Zag product is another classic example. But, of course, $\delta_1$ has some effect on the resulted code. The effect $\delta_1$ has on the code is via the query complexity. Indeed, as the analysis will show, the smaller $\delta_1$ is (i.e., the weaker the guarantee we have on the distance of $C_1$), the larger $k_2=k_2(\delta_1)$ and $k_3=k_3(\delta_1)$ must be, with a far stronger effect on $k_3$. More quantitatively, roughly speaking, by taking a sufficiently good sampler (e.g., the one that is given by Theorem 15), $k_2\approx\operatorname{poly}\log(1/\delta_1)$ and $k_3\approx\operatorname{poly}(1/\delta)$. This, in turn, effects the query complexities $q_2=q_2(k_2)$ and $q_3=q_3(k_3)$.
\begin{proof}[Author proof of Proposition 65]
That the query complexity is $q\leq q_1q_2q_3$ readily follows by the querying function, where recall that per our convention $q_3=q_3(2D)$. To analyze the rate, recall that $\rho_3$ is a non-decreasing function. Further, our convention dictates that by writing $\rho_3$ without explicitly mentioning the message length, we refer to $\rho_3$ applied with the smallest message length taken by the construction, namely, $\rho_3=\rho_3(D/2)$. Thus,
\[n=\sum_{v\in R}n_v=\sum_{v\in R}\frac{k_v}{\rho_3(k_v)}\leq\frac1{\rho_3}\sum_{v\in R}k_v=\frac{\ell n_2}{\rho_3}=\frac{n_1}{\rho_2\rho_3}.\]
Recall that $k=k_1=\rho_1n_1$ which shows that $\rho=k/n\geq\rho_1\rho_2\rho_3$.
We turn to analyze the distance $\delta$ and error $\varepsilon$. Let $x\in\Sigma^k$ and let $\widetilde{C(x)}\in\Sigma^n$ be such that $\operatorname{dist}(C(x),\widetilde{C(x)})\leq\delta$. Define the set of ``errors'', namely, the disagreements between $C(x)$ and $\widetilde{C(x)}$ by
\[B=\{i\in[n]\mid\widetilde{C(x)}_i\ne C(x)_i\}.\]
By assumption, $\mu(B)\leq\delta$. The error set $B$ induces errors ``backwards'' throughout the construction. We proceed by analyzing these induced errors. Recall that, in the encoding function, we defined for each $v\in[r]$ an element $t^{(v)}=t^{(v)}(x)\in\Sigma^{n_v}$. Partition $\widetilde{C(x)}$ to $r$ substrings $\widetilde{C(x)}=\widetilde t^{(1)}\circ\cdots\circ\widetilde t^{(r)}$, where $\widetilde t^{(v)}$ has length $n_v$, and define the set
\[B_t=\{v\in[r]\mid\operatorname{dist}(t^{(v)},\widetilde t^{(v)})\geq\delta_3\}.\]
Informally, $v\in B_t$ if the adversary has introduced too many errors on the respective block to allow for correct decoding via $D_3$.
\begin{ntclaim}[Author Claim 66] $\mu(B_t)\leq8\delta/\delta_3$.\end{ntclaim}
\begin{proof}
For $v\in R$ let $e_v=\operatorname{dist}(t^{(v)},\widetilde t^{(v)})$. We have that $\sum_{v\in R}e_vn_v\leq\delta n$. On the other hand,
\[\sum_{v\in R}e_vn_v\geq\delta_3\sum_{v\in B_t}n_v\geq\frac{\delta_3D|B_t|}{2}.\]
But, per our assumption that $\rho_3\geq1/2$, and since $k_v\leq2D$ for all $v\in R$,
\[n=\sum_{v\in R}n_v\leq2\sum_{v\in R}k_v\leq4Dr.\]
The proof follows by the above three inequalities.
\end{proof}
For convenience we also denote $B_w=B_t$. Next, we define
\begin{equation}B_z=\{u\in[\ell]\mid|\Gamma(u)\cap B_w|\geq\delta_2n_2\}.\tag{6.1}\end{equation}
\begin{ntclaim}[Author Claim 67] $\mu(B_z)\leq\delta_1$.\end{ntclaim}
\begin{proof}
By Claim 66 and by our assumption on $\delta$,
\[\mu(B_w)\leq\frac{8\delta}{\delta_3}=\frac{\delta_2}{2}.\]
Recall that $G$ is a $(\delta_2/2,\delta_1)$-sampler. Thus, at most $\delta_1$-fraction of the left vertices $u\in[\ell]$ satisfy $\mu(\Gamma(u)\cap B_w)\geq\mu(B_w)+\delta_2/2$. The proof then follows since $\mu(B_w)\leq\delta_2/2$.
\end{proof}
Lastly, define
\[B_y=\left\{a\in[n_1]\;\middle|\;\left\lceil\frac a{k_2}\right\rceil\in B_z\right\}.\]
For $v\in[r]$, $b\in[k_v]$ we define the function $\widetilde w_b^{(v)}:\{0,1\}^{r_3}\to\Sigma$ as follows: on input $s_3\in\{0,1\}^{r_3}$
\[\widetilde w_b^{(v)}(s_3)=D_3(b,\widetilde t^{(v)},s_3).\]
\begin{ntclaim}[Author Claim 68]
There exists a set $E_3\subseteq\{0,1\}^{r_3}$ with $\mu(E_3)\leq\varepsilon_3n$ such that for every $s_3\in\{0,1\}^{r_3}\setminus E_3$, $v\in[r]\setminus B_t$, and $b\in[k_3]$ it holds that $\widetilde w_b^{(v)}(s_3)=w_b^{(v)}$.
\end{ntclaim}
\begin{proof}
Fix $v\in[r]\setminus B_t$. By the definition of $B_t$, one has that $\operatorname{dist}(t^{(v)},\widetilde t^{(v)})\leq\delta_3$. By the encoding function, $t^{(v)}=C_3(w^{(v)})$. Therefore, for every $b\in[k_3]$,
\[\Pr_{s_3\sim U_{r_3}}[D_3(b,\widetilde t^{(v)},s_3)\ne w_b^{(v)}]\leq\varepsilon_3.\]
The proof then follows by taking the union bound over all $v\in[r]\setminus B_t$ and $b\in[k_v]$ as indeed $\sum k_v\leq n$.
\end{proof}
For $(u,j)\in[\ell]\times[n_2]$ we define the function $\widetilde z_j^{(u)}:\{0,1\}^{r_3}\to\Sigma$ as follows. For $s_3\in\{0,1\}^{r_3}$ we have $\widetilde z_j^{(u)}(s_3)=\widetilde w_{j'}^{(v)}(s_3)$, where $(v,j')=\Gamma(u,j)\in[r]\times[k_v]$. Further define the function $\widetilde z^{(u)}:\{0,1\}^{r_3}\to\Sigma^{n_2}$ by
\[\widetilde z^{(u)}(s_3)=\widetilde z_1^{(u)}(s_3)\circ\cdots\circ\widetilde z_{n_2}^{(u)}(s_3).\]
\begin{ntclaim}[Author Claim 69]
For every $u\notin B_z$ and $s_3\in\{0,1\}^{r_3}\setminus E_3$ it holds that
\[\operatorname{dist}(\widetilde z^{(u)}(s_3),z^{(u)})\leq\delta_2.\]
\end{ntclaim}
\begin{proof}
Fix $s_3\in\{0,1\}^{r_3}\setminus E_3$ and consider any $u\in[\ell]\setminus B_z$. By the encoding function, for every $j\in[n_2]$ it holds that $z_j^{(u)}=w_{j'}^{(v)}$, where $(v,j')=\Gamma(u,j)$. As $v\notin B_z$, at most $\delta_2n_2$ of $j\in[n_2]$ satisfy $v\in B_w$. For every other $j$,
\[\widetilde z_j^{(u)}=\widetilde w_{j'}^{(v)}(s_3)=w_{j'}^{(v)}=z_j^{(u)},\]
proving the claim.
\end{proof}
For $u\in[\ell]$, $a\in[k_2]$ we define the function $\widetilde y_a^{(u)}:\{0,1\}^{r_2}\times\{0,1\}^{r_3}\to\Sigma$ as follows. On $(s_2,s_3)\in\{0,1\}^{r_2}\times\{0,1\}^{r_3}$,
\[\widetilde y_a^{(u)}(s_2,s_3)=D_2(a,\widetilde z^{(u)}(s_3),s_2).\]
\begin{ntclaim}[Author Claim 70]
There exists a set $E_2\subseteq\{0,1\}^{r_2}$ with $\mu(E_2)\leq\varepsilon_2n$ such that for every $u\in[\ell]\setminus B_z$, $a\in[k_2]$, and $(s_2,s_3)\in(\{0,1\}^{r_2}\setminus E_2)\times(\{0,1\}^{r_3}\setminus E_3)$ it holds that $\widetilde y_a^{(u)}(s_2,s_3)=y_a^{(u)}$.
\end{ntclaim}
\begin{proof}
Fix $u\in[\ell]\setminus B_z$. By the encoding function $z^{(u)}=C_2(y^{(u)})$. Recall that
\[\widetilde y_a^{(u)}(s_2,s_3)=D_2(a,\widetilde z^{(u)}(s_3),s_2).\]
As $s_3\notin E_3$, $u\notin B_z$, Claim 69 implies $\operatorname{dist}(\widetilde z^{(u)}(s_3),z^{(u)})\leq\delta_2$. Therefore
\[\Pr_{s_2\sim U_{r_2}}[D_2(a,\widetilde z^{(u)}(s_3),s_2)\ne y_a^{(u)}]\leq\varepsilon_2.\]
The proof then follows by taking the union bound over all $a\in[k_2]$ and $u\in[\ell]\setminus B_z$, and noting that $k_2\ell=n_1\leq n$.
\end{proof}
\begin{ntclaim}[Author Claim 71]
For every $(s_2,s_3)\in(\{0,1\}^{r_2}\setminus E_2)\times(\{0,1\}^{r_3}\setminus E_3)$, it holds that
\[\operatorname{dist}(\widetilde y(s_2,s_3),y)\leq\delta_1,\]
where $\widetilde y(s_2,s_3)$ is the concatenation of the $k_2$-length strings $(\widetilde y^{(u)}(s_2,s_3)\mid u\in[\ell])$.
\end{ntclaim}
\begin{proof}
Note that by Claim 70, the projection of the two strings $\widetilde y(s_2,s_3),y$ to a block corresponding to $u\notin B_z$ are in full agreement. The proof then follows by Claim 67.
\end{proof}
We now conclude the proof of Proposition 65. Let $p\in[k]$, by Claim 71, for every $(s_2,s_3)\in(\{0,1\}^{r_2}\setminus E_2)\times(\{0,1\}^{r_3}\setminus E_3)$, we have that $\operatorname{dist}(\widetilde y(s_2,s_3),y)\leq\delta_1$. Since by the encoding function $y=C_1(x)$, it holds
\[\Pr_{s_1\sim U_{r_1}}[D_1(p,\widetilde y(s_2,s_3),s_1)\ne x_p]\leq\varepsilon_1.\]
The proof then follows since $\mu(E_2)\leq\varepsilon_2n$ and $\mu(E_3)\leq\varepsilon_3n$.
\end{proof}

\subsection*{6.2.1 Proof of Theorem 5}
In this short section we prove Theorem 5. We focus on the version that is based on non-explicit samplers, yielding non-explicit reductions. The explicit reduction, which entails a bit more technical work, is deferred to Section 6.3 and Section 6.6. We choose to focus on the non-explicit version first because we believe that understanding LDC in the information-theoretic level is, at present, a deeper and more urgent problem than the question of explicitness. Also, the parameters are easier to work with. For the information-theoretic version, we make use of the sampler that is given by Theorem 15. From here on we refer to the constant $c_{\mathrm{samp}}\geq1$ that appears in that theorem.
\begin{proof}[Author proof of Theorem 72, stated at the start of this item]
Take $C_1$ to be the code $C$ in the hypothesis of the theorem, namely, a code with block length $n_1=n$ and distance $\delta_1=\delta$. Recall that in the distance amplification procedure from Section 6.1, we make use of a $(\delta_2/2,\delta_1)$ sampler $G=([\ell],[r],E)$ with $\ell=n_1/k_2$ and left-degree $d=n_2$. For the proof, we will instantiate the distance amplification procedure with the sampler that is given by Theorem 15. We take $C_2$ to be an asymptotically good code over $\Sigma$ set with block length
\[n_2=c_{\mathrm{samp}}\cdot\frac{\log(1/\delta_1)}{(\delta_2/2)^2}=O(\log(1/\delta)),\]
where $\delta_2$ is the (constant) distance of $C_2$, having rate at least $1/2$. Note that this choice of parameters is as required by Theorem 15 from the left degree of the sampler. Clearly, $C_2$ has query complexity $O(\log(1/\delta))$ and error $\varepsilon_2=0$. As for the degree $D_v$ of any given right vertex $v$ of the sampler, note that the average right degree is
\[D=\frac{\ell d}{r}=\frac d{\delta\log(1/\delta)}=\frac{4c_{\mathrm{samp}}}{\delta_1\delta_2^2}=\Theta\left(\frac1\delta\right).\]
Recall that, by Theorem 15, $D_v\in[D/2,2D]$. For every length in this range, we take a code from the family $C'$ having the required message length. We would like take the family of codes $C_3$ to be $C'$ though we must reduce the error first. Indeed, note that the error $\varepsilon$ of the code obtained by Proposition 65 is $\varepsilon_1+n(\varepsilon_2+\varepsilon_3)$. As mentioned in the introduction, one can reduce the error from $1/5$ to $1/(10n)$ by applying the decoding procedure for $c\log n$ times, where $c$ is some large enough constant, and output the symbol according to plurality. This has no effect on the rate or distance of $C'$, and has a multiplicative $O(\log n)$ cost in query complexity. That is, the query complexity of $C_3$ is $O(q'_{O(1/\delta_1)}\log n)$. The proof then readily follows by Proposition 65.
\end{proof}
\subsubsection*{6.2.1.1 Improving the query complexity further given low-error LDC}
We remark that, if $C'$ has error $O(1/n)$ to begin with, $n$ being the block length of $C$, then one can skip the error reduction in the proof of Theorem 72, and get a slightly better query complexity. Indeed, this will save the $\log n$ factor in Equation (6.2). Moreover, observe that $C_2$ can be taken to be an LDC as well, rather than a standard code, which will reduce its deterioration on the query complexity from $O(\log(1/\delta))$ to $q'_{O(\log(1/\delta))}$. However, for that, one need the error of $C_2$ to be $O(1/n)$ as well. Assuming one can obtain such low-error LDC (note that an error of $1/n$ is at least exponentially-small in the length of $C_2$ since $\delta>1/n$), the query complexity can be improved further to
\[q_{\mathrm{new}}\leq q\cdot q'_{O(1/\delta)}q'_{O(\log(1/\delta))}.\]
We conclude this section by instantiating Theorem 72 with $C'$ taken to be the state-of-the-art construction of asymptotically good LDC.
\setcounter{nttheorem}{72}
\begin{nttheorem}[Author Theorem 73, {\cite{ref26}}]
Let $\Sigma$ be a finite alphabet. Then, there exist constants $\delta,\rho$ and an explicit infinite family of $(q_k,\delta,1/5)$-LDC, $k$ being the message length, having query complexity $q_k=2^{O(\sqrt{\log(k)\log\log k})}$.
\end{nttheorem}
Using it, one gets query complexity
\[q_{\mathrm{new}}\leq q\log(n)\cdot2^{O(\sqrt{\log(1/\delta)\cdot\log\log(1/\delta)})}=q\log(n)(1/\delta)^{o(1)}.\]

\begin{thebibliography}{43}
\bibitem{ref1}Noga Alon, Jeff Edmonds, and Michael Luby. Linear time erasure codes with nearly optimal recovery. In Proceedings of IEEE 36th Annual Foundations of Computer Science, pages 512–519. IEEE, 1995.
\bibitem{ref2}Noga Alon and Michael Luby. A linear time erasure-resilient code with nearly optimal recovery. IEEE Transactions on Information Theory, 42(6):1732–1736, 1996.
\bibitem{ref3}Sanjeev Arora, Carsten Lund, Rajeev Motwani, Madhu Sudan, and Mario Szegedy. Proof verification and the hardness of approximation problems. Journal of the ACM (JACM), 45(3):501–555, 1998.
\bibitem{ref4}Sanjeev Arora and Shmuel Safra. Probabilistic checking of proofs: A new characterization of NP. Journal of the ACM (JACM), 45(1):70–122, 1998.
\bibitem{ref5}László Babai, Lance Fortnow, Leonid A Levin, and Mario Szegedy. Checking computations in polylogarithmic time. In Proceedings of the twenty-third annual ACM symposium on Theory of computing, pages 21–32, 1991.
\bibitem{ref6}Laszlo Babai, Lance Fortnow, Noam Nisan, and Avi Wigderson. BPP has subexponential time simulations unless EXPTIME has publishable proofs. Computational Complexity, 3(4):307–318, 1993.
\bibitem{ref7}Omer Barkol, Yuval Ishai, and Ronny Roth. Locally decodable codes and their applications. PhD thesis, Computer Science Department, Technion, 2008.
\bibitem{ref8}Donald Beaver and Joan Feigenbaum. Hiding instances in multioracle queries. In Annual Symposium on Theoretical Aspects of Computer Science, pages 37–48. Springer, 1990.
\bibitem{ref9}Mihir Bellare and John Rompel. Randomness-efficient oblivious sampling. In Proceedings of the 35th Annual Symposium on Foundations of Computer Science, 1994, pages 276–287. IEEE, 1994.
\bibitem{ref10}Manuel Blum, Michael Luby, and Ronitt Rubinfeld. Self-testing/correcting with applications to numerical problems. In Proceedings of the twenty-second annual ACM symposium on Theory of computing, pages 73–83, 1990.
\bibitem{ref11}Benny Chor, Oded Goldreich, Eyal Kushilevitz, and Madhu Sudan. Private information retrieval. In Proceedings of IEEE 36th Annual Foundations of Computer Science, pages 41–50. IEEE, 1995.
\bibitem{ref12}Zeev Dvir. On matrix rigidity and locally self-correctable codes. computational complexity, 20(2):367–388, 2011.
\bibitem{ref13}Zeev Dvir, Parikshit Gopalan, and Sergey Yekhanin. Matching vector codes. SIAM Journal on Computing, 40(4):1154–1178, 2011.
\bibitem{ref14}Klim Efremenko. 3-query locally decodable codes of subexponential length. SIAM Journal on Computing, 41(6):1694–1703, 2012.
\bibitem{ref15}Peter Gemmell, Richard Lipton, Ronitt Rubinfeld, Madhu Sudan, and Avi Wigderson. Self- testing/correcting for polynomials and for approximate functions. In STOC, volume 91, pages 32–42. Citeseer, 1991.
\bibitem{ref16}Peter Gemmell and Madhu Sudan. Highly resilient correctors for polynomials. Information processing letters, 43(4):169–174, 1992.
\bibitem{ref17}Oded Goldreich. A sample of samplers: A computational perspective on sampling. In Studies in Complexity and Cryptography, pages 302–332. Springer, 2011. doi:10.1007/ 978-3-642-22670-0\_24.
\bibitem{ref18}Oded Goldreich and Madhu Sudan. Locally testable codes and PCPs of almost-linear length. Journal of the ACM (JACM), 53(4):558–655, 2006.
\bibitem{ref19}Sivakanth Gopi, Swastik Kopparty, Rafael Oliveira, Noga Ron-Zewi, and Shubhangi Saraf. Locally testable and locally correctable codes approaching the gilbert-varshamov bound. IEEE Transactions on Information Theory, 64(8):5813–5831, 2018.
\bibitem{ref20}Alan Guo, Swastik Kopparty, and Madhu Sudan. New affine-invariant codes from lifting. In Proceedings of the 4th conference on Innovations in Theoretical Computer Science, pages 529–540. ACM, 2013.
\bibitem{ref21}Venkatesan Guruswami, Atri Rudra, and Madhu Sudan. Essential coding theory. Draft available at http://www.cse.buffalo.edu/\textasciitilde{}atri/courses/coding-theory/book, 2012.
\bibitem{ref22}Brett Hemenway, Rafail Ostrovsky, and Mary Wootters. Local correctability of expander codes. Information and Computation, 243:178–190, 2015.
\bibitem{ref23}Jonathan Katz and Luca Trevisan. On the efficiency of local decoding procedures for error- correcting codes. In Proceedings of the thirty-second annual ACM symposium on Theory of computing, pages 80–86, 2000.
\bibitem{ref24}Tali Kaufman and Madhu Sudan. Sparse random linear codes are locally decodable and testable. In 48th Annual IEEE Symposium on Foundations of Computer Science (FOCS’07), pages 590–600. IEEE, 2007.
\bibitem{ref25}Kiran S Kedlaya and Sergey Yekhanin. Locally decodable codes from nice subsets of finite fields and prime factors of mersenne numbers. SIAM Journal on Computing, 38(5):1952–1969, 2009.
\bibitem{ref26}Swastik Kopparty, Or Meir, Noga Ron-Zewi, and Shubhangi Saraf. High-rate locally correctable and locally testable codes with sub-polynomial query complexity. Journal of the ACM (JACM), 64(2):11, 2017.
\bibitem{ref27}Swastik Kopparty, Shubhangi Saraf, and Sergey Yekhanin. High-rate codes with sublinear-time decoding. Journal of the ACM (JACM), 61(5):28, 2014.
\bibitem{ref28}Ray Li and Mary Wootters. Lifted multiplicity codes and the disjoint repair group property. In Approximation, Randomization, and Combinatorial Optimization. Algorithms and Techniques (APPROX/RANDOM 2019). Schloss Dagstuhl-Leibniz-Zentrum fuer Informatik, 2019.
\bibitem{ref29}Richard J. Lipton. Efficient checking of computations. In Annual Symposium on Theoretical Aspects of Computer Science, pages 207–215. Springer, 1990.
\bibitem{ref30}Irving S Reed. A class of multiple-error-correcting codes and the decoding scheme. Technical report, Massachusetts inst of tech Lexington Lincoln lab, 1953.
\bibitem{ref31}Omer Reingold, Salil Vadhan, and Avi Wigderson. Entropy waves, the zig-zag graph product, and new constant-degree expanders and extractors. In Electronic Colloquium on Computational Complexity (ECCC), page 18, 2001. URL: https://eccc.weizmann.ac.il/report/2001/018/.
\bibitem{ref32}Ronitt Rubinfeld and Madhu Sudan. Robust characterizations of polynomials with applications to program testing. SIAM Journal on Computing, 25(2):252–271, 1996.
\bibitem{ref33}C. E. Shannon. A mathematical theory of communication. ACM SIGMOBILE Mobile Computing and Communications Review, 5(1):3–55, 2001. Originally appeared in Bell System Tech. J. 27:379–423, 623–656, 1948.
\bibitem{ref34}Carl Siegel. Über die classenzahl quadratischer zahlkörper. Acta Arithmetica, 1(1):83–86, 1935.
\bibitem{ref35}Madhu Sudan, Luca Trevisan, and Salil Vadhan. Pseudorandom generators without the XOR lemma. Journal of Computer and System Sciences, 62(2):236–266, 2001.
\bibitem{ref36}Luca Trevisan. List-decoding using the XOR lemma. In 44th Annual IEEE Symposium on Foundations of Computer Science, 2003. Proceedings., pages 126–135. IEEE, 2003.
\bibitem{ref37}Arnold Walfisz. Zur additiven zahlentheorie. ii. Mathematische Zeitschrift, 40(1):592–607, 1936.
\bibitem{ref38}David Woodruff. New lower bounds for general locally decodable codes. In Electronic Colloquium on Computational Complexity (ECCC), volume 14, 2007.
\bibitem{ref39}S. Yekhanin. Locally decodable codes. In International Computer Science Symposium in Russia, pages 289–290. Springer, 2011.
\bibitem{ref40}Sergey Yekhanin. Towards 3-query locally decodable codes of subexponential length. Journal of the ACM (JACM), 55(1):1–16, 2008.
\bibitem{ref41}Kalina Petrova Zeev Dvir. Lecture 1: Introduction. Lecture notes: https://www.cs. princeton.edu/\textasciitilde{}zdvir/LDCnotes/LDC1.pdf, year = 2016.
\bibitem{ref42}Kalina Petrova Zeev Dvir. Lecture 4: Lower bounds for r-query LDCs. Lecture notes: https://www.cs.princeton.edu/\textasciitilde{}zdvir/LDCnotes/LDC4.pdf, 2016.
\bibitem{ref43}Victor Vasilievich Zyablov. An estimate of the complexity of constructing binary linear cascade codes. Problemy Peredachi Informatsii, 7(1):5–13, 1971.
\end{thebibliography}
\end{document}
